\documentclass[12pt,a4paper]{article}
\usepackage[utf8]{inputenc}
\usepackage[T1]{fontenc}
\usepackage[french]{babel}
\usepackage{amsmath}
\usepackage{geometry}
\usepackage{setspace}
\usepackage{hyperref}
\geometry{margin=2.5cm}
\setstretch{1.25}

\title{Revue de Synthèse : Mobilité, salaires et trajectoires professionnelles}
\author{Thomas, Chloé \& Ilayda}
\date{Octobre 2025}

\begin{document}
\maketitle


\section{Les approches de Glaude (1986) et Simonnet (1996) : capital humain, ancienneté et mobilité}

Les travaux de \textbf{Glaude (1986)} et de \textbf{Simonnet (1996)} s’inscrivent dans la continuité de la théorie du capital humain développée par Becker (1975) et Mincer (1974).  
Tous deux cherchent à expliquer la formation des salaires à partir des caractéristiques individuelles des travailleurs, notamment la formation, l’expérience et l’ancienneté.  
Là où Glaude met en évidence la rentabilité de la stabilité et du capital humain spécifique, Simonnet prolonge cette réflexion en y intégrant la mobilité professionnelle et les différences entre hommes et femmes.

\subsection{Le modèle de Glaude (1986)}

Dans son article \textit{L’ancienneté et l’expérience professionnelle dans la formation du salaire}, Glaude propose une estimation économétrique visant à distinguer les effets de l’expérience professionnelle totale, qui reflète le \textbf{capital humain général}, et ceux de l’ancienneté dans l’entreprise, qui mesurent le \textbf{capital humain spécifique}.  
Son modèle de base, inspiré de la fonction de gains de Mincer, prend la forme suivante :

\begin{equation}
\ln W_i = \alpha + \beta_1 \text{Exp}_i + \beta_2 \text{Exp}_i^2 + \gamma_1 \text{Anc}_i + \gamma_2 \text{Anc}_i^2 + \delta' X_i + \varepsilon_i
\end{equation}

où :
\begin{itemize}
    \item $W_i$ est le salaire horaire de l’individu $i$ ;
    \item $\text{Exp}_i$ représente l’expérience professionnelle totale (capital humain général) ;
    \item $\text{Anc}_i$ désigne l’ancienneté dans l’entreprise (capital humain spécifique) ;
    \item $X_i$ regroupe les variables de contrôle : diplôme, catégorie professionnelle, secteur, taille d’entreprise et région ;
    \item les termes au carré ($\text{Exp}_i^2$ et $\text{Anc}_i^2$) traduisent la décroissance des rendements avec la durée de carrière.
\end{itemize}

\ 

Les résultats indiquent que l’expérience et l’ancienneté ont toutes deux un effet positif et significatif sur les salaires, mais de manière différenciée.  
Le rendement de l’expérience est estimé à environ \textbf{1,5 à 2\% par an}, tandis que celui de l’ancienneté se situe autour de \textbf{1\%}.  
Les coefficients quadratiques sont négatifs, ce qui confirme la présence de rendements décroissants à mesure que la durée dans l’emploi augmente.  

Glaude montre également que ces rendements varient selon les secteurs d’activité : ils sont plus élevés dans les \textbf{secteurs fermés ou hiérarchisés} (banques, administration) que dans les secteurs concurrentiels (industrie, BTP).  
Cependant, il met en garde contre une possible \textit{surestimation} du rendement de l’ancienneté : une longue présence dans une entreprise peut refléter un apprentissage productif, mais aussi une bonne adéquation entre le salarié et son emploi.  
Cette observation ouvre la voie aux travaux ultérieurs de Simonnet.

\subsection{Le prolongement de Simonnet (1996)}

Dix ans plus tard, \textbf{Simonnet (1996)} reprend cette distinction entre expérience générale et capital humain spécifique, mais elle y ajoute une dimension absente chez Glaude : celle du \textbf{déroulement de carrière} et de la \textbf{mobilité professionnelle}.  
Elle part du même cadre de Mincer, en y intégrant explicitement les mouvements de mobilité interne et externe.  
Son modèle théorique s’écrit ainsi :

\begin{equation}
\ln E_{it} = \ln E_{i0} + r_s S_i + r_e Ex_i + r_a Anc_i + \eta M_i + u_{it}
\end{equation}

où :
\begin{itemize}
    \item $S_i$ désigne le nombre d’années d’études de l’individu $i$ ;
    \item $Ex_i$ et $Anc_i$ représentent respectivement l’expérience totale et l’ancienneté ;
    \item $M_i$ correspond au profil de mobilité (interne, externe, mixte ou public/privé).
\end{itemize}

\

Simonnet montre que si l’on ne prend pas en compte la mobilité, on risque de mal interpréter le rôle de l’ancienneté.
Un salarié qui reste longtemps dans la même entreprise peut être mieux payé, non seulement parce qu’il a accumulé de l’expérience spécifique, mais aussi parce qu’il occupe un poste bien adapté à ses compétences. Ses résultats confirment cette idée.

\ 

Chez les \textbf{hommes}, l’effet de l’ancienneté diminue fortement dès qu’on ajoute les variables de mobilité : la hausse des salaires provient surtout de la \textbf{mobilité interne}, liée aux promotions ou à la progression hiérarchique.
Chez les \textbf{femmes}, c’est plutôt la \textbf{mobilité externe}, le changement d’entreprise qui permet d’obtenir les meilleurs gains salariaux.
Enfin, passer du \textbf{secteur public vers le privé} est en général plus avantageux que l’inverse.

\subsection{Une continuité analytique entre Glaude et Simonnet}

Les travaux de Glaude et de Simonnet sont liés : tous deux étudient comment l’expérience et la stabilité influencent les salaires.
Simonnet reprend le modèle de Glaude, mais elle l’enrichit en ajoutant la mobilité professionnelle et la différence entre hommes et femmes.
Les deux montrent que les salaires augmentent avec l’expérience et l’ancienneté, mais que ces effets dépendent du parcours de chaque individu.
Leurs résultats diffèrent sur le rôle de l’ancienneté.
Pour Glaude, elle reflète surtout la valeur de la stabilité dans l’entreprise.
Pour Simonnet, cet effet devient plus faible quand on prend en compte la mobilité : les salariés les mieux payés sont souvent ceux qui ont trouvé un emploi bien adapté à leurs compétences et qui y restent plus longtemps. 

\subsection{Conclusion}
 
Glaude montre que rester longtemps dans la même entreprise peut être un bon investissement : la stabilité et la fidélité sont récompensées.
Simonnet, elle, apporte une vision plus dynamique. Elle explique que la mobilité joue aussi un rôle important pour comprendre les différences de salaires entre hommes et femmes.
Ainsi, Simonnet ne contredit pas Glaude, mais elle complète son analyse.
Elle montre que la stabilité n’est pas toujours rentable : tout dépend du type de mobilité, du secteur et du genre.
L’ancienneté reste importante, mais son effet sur le salaire doit être étudié en tenant compte du parcours professionnel et des possibilités de changement d’emploi.



\end{document}